% !TEX root = ../notes_missingsemester.tex
% Chapter: HighFive — Contribute to a Real Full Stack Application
%%%%%%%%%%%%%%%%%%%%%%%%%%%%%%%%%%%%%%%%%%%%%%%%%%%%%%%%%%%%%%%%%%%%%%

\index{HighFive}\index{citizen science}\index{contribution}
\chapter{Ship on something that runs}
\printchapterversion{11_highfive_contribution}
\newthought{High Five.}

\section{The Why}\index{citizen science}\index{wild bees}

The Projektbericht lets you extend a block of this course on your own snapshot.
HighFive offers the other path: A real full stack application, running live at
\href{https://highfive.schutera.com}{highfive.schutera.com}, with its source on my
GitHub and an open invitation to ship into it. Where the report proves you can build
in isolation, a contribution proves you can build inside a system that other people
already depend on.

\marginnote{\newthought{Find HighFive here.} The deployment runs at
\href{https://highfive.schutera.com}{highfive.schutera.com}, and the source lives at
\href{https://github.com/schutera/highfive}{github.com/schutera/highfive}.

\medskip
\null\hfill\qrcode[height=1.0cm]{https://github.com/schutera/highfive}
\medskip

Read the arc42 architecture document in the repository before you touch code; 
it explains how the pieces fit together.}

\newthought{HighFive is an open monitoring pipeline} for wild bee nesting activity,
built as a citizen science kit. Camera modules running on solar power photograph bee hotels,
an analysis pipeline turns those images into counts and trends, and an interactive dashboard
with a map puts the results in front of anyone who cares about pollinators. 

\newthought{It is the full stack you just learned}, end to end.
The frontend is React with Vite and TypeScript. The backend is Node.js with Express
alongside Python services on Flask. Data lives in DuckDB. Firmware for the ESP32-CAM
camera modules is written in C++. Every layer from this course, from the Linux box up to the
served frontend, has a home somewhere in this repository.

\marginnote{\newthought{Before you open a pull request.} HighFive is source available rather
than open source: Free for any noncommercial use under the PolyForm Noncommercial License
1.0.0, with commercial use under a separate license. Contributions require you to accept the
Contributor License Agreement in the repository. For commercial questions, write to
\href{mailto:schutera@dhbw-ravensburg.de}{schutera@dhbw-ravensburg.de}.}

\newthought{Pick something and make it better.}
Sharpen the image pipeline, harden a backend service, tidy the dashboard, document a rough edge,
or write the test that was missing. Open an issue describing what you want to do,
agree the scope with me, then send a pull request. By prior agreement a merged contribution
can form the basis of your Projektbericht: Same evaluation, real impact, and your name in the
history of a project that outlives the semester. And yes, other students 
did ship into this repo in previous semesters, so you will be in good company.
